\documentclass[10pt,a4paper]{amsart}
\usepackage[margin=19mm]{geometry}
\usepackage{amsmath,amssymb,mathtools}
\usepackage[scr=rsfs,cal=euler]{mathalfa}
\usepackage{xfrac}
\usepackage[unicode,hidelinks,bookmarks=false]{hyperref}
\newcommand{\ie}{i.e.\ }
\newcommand{\cf}{cf.\ }
\newcommand{\resp}{respectively\ }
\newcommand{\Subsection}{\subsection}
\providecommand{\nobreakdash}{\nobreak\hbox{-}}
\setlength{\emergencystretch}{2em}
\multlinegap0pt
\usepackage{amssymb}
\usepackage{enumerate}
\newtheorem{lemma}{Lemma}
\title{Explicit bounds for Dickman's function}
\author{Andreas Weingartner}
\date{Source: arXiv:2606.07785v3 (CC BY 4.0)}
\begin{document}
\maketitle

\noindent\emph{Source and licence.} Explicit bounds for Dickman's function. Version arXiv:2606.07785v3 (CC BY 4.0), distributed by arXiv under the Creative Commons Attribution 4.0 International licence, http://creativecommons.org/licenses/by/4.0/, as recorded in the arXiv licence metadata for this exact version and shown on the abstract page https://arxiv.org/abs/2606.07785v3. Full licence notice: \texttt{licenses/arxiv-2606.07785-CC-BY-4.0.txt}.\par\smallskip

\noindent\textit{Editorial scope.} Counted unit: the article's lemma lemalpha (source0.tex lines 504--511). The article's complete original proof of it is delivered with it (source0.tex lines 512--531, one \texttt{proof} environment of 20 lines). No mathematics is written, repaired or re-derived: the statement and the proof are the article's own lines, in the article's own notation, and no retained line is substituted or re-flowed.  Retained uncounted as the source-local context the proof needs: lines 196--207; lines 213--218; lines 481--502.  Nothing was omitted from any retained span, so the package carries no omission record.  No retained line contains a citation, so the wrapper carries no bibliography and no bibliography entry.  No retained line contains a figure environment, an image-inclusion command or drawing code, so no image is delivered and the package declares no asset dependency.  Editorial formatting only, in addition: an AMS \texttt{amsart} preamble that restates the article's own declarations and macros used by the retained lines, namely the environments \texttt{lemma} on separate counters, together with the standard packages those lines need.  The retained environments print in this document as Lemma 1 in order of appearance, whereas the article prints the same items with the numbering of its own full document.  The complete native source of the article as distributed in the arXiv e-print for this exact version is bundled unchanged as \texttt{sources/202274/source0.tex}.
\begin{lemma}\label{lemIk}
For real $x>0$ and natural numbers $k$ we have 
$$
\frac{e^x -1}{x} \ge I^{(k)}(x) \ge I^{(k+1)}(x) \quad (k\ge 1),
$$
$$
I^{(k)}(x) \ge \frac{e^x}{x} \left(1-\frac{k-1}{x}\right) \qquad (k\ge 2),
$$
$$
I^{(k)}(x) \le  \frac{e^x}{x} \left(1-\frac{k-1}{x}+\frac{(k-1)(k-2)}{x^2}\right) \qquad (k\ge 3).
$$
\end{lemma}

\begin{lemma}\label{lemxi}
For $u\ge 2$ we have 
$$
\log (u\log u) < \xi(u) < \log (u\log u) +1.
$$
\end{lemma}

\begin{lemma}\label{lemAx}
For $u\ge 2$ we have 
$$
 -\frac{\xi}{12 e^\xi}\left(1-\frac{1}{\xi}+\frac{1.75}{\xi^2}\right)< \alpha(u)< -\frac{\xi}{12 e^\xi}\left(1-\frac{1}{\xi}+\frac{0.5}{\xi^2}\right) .
$$
\begin{proof}
Let
$$
A(x):=\frac{1}{8} \frac{I^{(4)}(x)}{(I^{(2)}(x))^2} - \frac{5}{24} \frac{(I^{(3)}(x))^2}{(I^{(2)}(x))^3} ,
$$
so that $\alpha(u)=A(\xi(u))$, and
define $\kappa(x)$ by 
$$
A(x) = -\frac{x}{12 e^x}\left(1-\frac{1}{x}+\frac{\kappa(x)}{x^2}\right).
$$
Lemma \ref{lemIk} shows that $\kappa(x)$ is bounded.
A calculus exercise, aided by a computer, reveals that 
$\kappa(x)$ is increasing on the interval $(0, 2.749...)$, where it reaches its maximum of $1.746...$ at $x=2.749...$, and decreasing 
on $(2.749..., \infty)$ with $\lim_{x\to \infty} \kappa(x) = 1/2$. We have $\kappa(1/3)=0.511...>0.5$. 
The result now follows since $\xi >1/3$ for $u\ge 2$. 
\end{proof}
\end{lemma}

\begin{lemma}\label{lemalpha}
For $u\ge 1$ we have 
$$
-\frac{1}{12 u}
<\alpha(u)-\frac{0.02}{u^2} < \alpha(u)+\frac{0.02}{u^2} 
< - \frac{1}{12 u (1+1/\log u)}.
$$
\end{lemma}

\begin{proof}
For $u\ge 4$, we have $u^2(1-1.75/\xi) >(0.02)  12(u\xi +1)=(0.02) 12  e^\xi$, by Lemma \ref{lemxi}. Lemma \ref{lemAx} 
implies
$$
\alpha(u) -\frac{0.02}{u^2}>  -\frac{\xi}{12 e^\xi}\left(1-\frac{1}{\xi}+\frac{1.75}{\xi^2}\right)-\frac{0.02}{u^2}
>  \frac{-\xi}{12 e^\xi}=\frac{-\xi}{12 (u\xi+1)}> \frac{-1}{12 u}.
$$
For $u\ge 1$, we have $u^2\ge (0.02) 24 \xi e^\xi= (0.02) 24 \xi(u\xi+1)$. Lemma \ref{lemAx} yields
$$
\alpha(u) +\frac{0.02}{u^2} < -\frac{\xi}{12 e^\xi}\left(1-\frac{1}{\xi}+\frac{0.5}{\xi^2}\right)+\frac{0.02}{u^2}
< \frac{-\xi}{12 e^\xi}\left(1-\frac{1}{\xi}\right)=\frac{-(\xi -1)}{12(u\xi+1)}.
$$
We claim that 
$$
-\frac{\xi -1}{12(u\xi+1)}< -\frac{1}{12u(1+1/\log u)}.
$$
Indeed, this is equivalent to $\xi>1+(1+1/u)\log u$, which holds for $u\ge 9$, by Lemma \ref{lemxi}.
This demonstrates the result for $u\ge 9$. For $1 \le u\le 9$ we confirm the result by graphing the
four quantities in question multiplied by $u$. 
\end{proof}

\end{document}
